%===============================================================================================================================%
%                                                            PREAMBLE                                                           %
%===============================================================================================================================%

\documentclass[honours,12pt]{unswthesis}
\linespread{1}
\usepackage{amsfonts}
\usepackage{amssymb}
\usepackage{amsthm}
\usepackage{latexsym,amsmath}
\usepackage{graphicx}
\usepackage{afterpage}
\usepackage{cite}
\usepackage{listings}
\usepackage{xcolor}
\usepackage[many]{tcolorbox}
\usepackage[colorlinks=true,linkcolor=blue,draft=false]{hyperref}
\usepackage[nameinlink]{cleveref}
\usepackage[ruled,vlined]{algorithm2e}
\usepackage{MnSymbol}

\hypersetup{citecolor=[rgb]{0, 0.5, 0}}

\crefname{algocfline}{algorrithm}{algorithms}
\Crefname{algocfline}{Algorithm}{Algorithms}

% Macros
\newcommand{\N}{\mathbb{N}} % Natural numbers
\renewcommand{\P}{\mathbb{P}} % Prime numbers
\newcommand{\Z}{\mathbb{Z}} % Integers
\newcommand{\Q}{\mathbb{Q}} % Rational numbers
\newcommand{\A}{\mathbb{A}} % Algebraic numbers
\newcommand{\B}{\mathbb{B}} % Algebraic integers
\newcommand{\R}{\mathbb{R}} % Real numbers
\newcommand{\C}{\mathbb{C}} % Complex numbers
\newcommand{\T}{\mathbb{T}} % Circle group
\newcommand{\F}{\mathbb{F}} % Field

\newcommand{\order}{\mathcal{O}} % Order

\newcommand{\GL}{\text{GL}} % General linear group
\newcommand{\SL}{\text{SL}} % Special linear group

\newcommand{\pmat}[1]{\begin{pmatrix}#1\end{pmatrix}}
\newcommand{\bmat}[1]{\begin{bmatrix}#1\end{bmatrix}}

\newcommand{\lnote}[1]{\llap{$#1$\hspace{0.75cm}}}
\newcommand{\rnote}[1]{\rlap{\hspace{0.75cm}$#1$}}

\newcommand{\red}[1]{\textcolor{red}{#1}}
\newcommand{\todo}[1]{\red{to do: #1}}
\newcommand{\tocheck}{(\red{check})}

\newcommand{\blankpage}{\null\thispagestyle{empty}\addtocounter{page}{-1}\newpage}


% Environments
\newtheorem{theorem}{Theorem}[section]
\newtheorem{lemma}[theorem]{Lemma}
\newtheorem{proposition}[theorem]{Proposition}
\newtheorem{corollary}[theorem]{Corollary}
\newtheorem{conjecture}[theorem]{Conjecture}
\newtheorem{definition}[theorem]{Definition}
\newtheorem{example}[theorem]{Example}
\newtheorem{remark}[theorem]{Remark}
\newtheorem{question}[theorem]{Question}
\newtheorem{notation}[theorem]{Notation}
\numberwithin{equation}{section}


%===============================================================================================================================%
%                                                            PREFACE                                                            %
%===============================================================================================================================%

% Title page
\title{Factoring via Binary Quadratic Forms}
\authornameonly{Harrison Graham}
\author{\Authornameonly\\{\bigskip}Supervisor: Professor David Harvey}
\copyrightfalse
\figurespagefalse
\tablespagefalse


% Plagiarism statement
\begin{document}
\beforepreface

\prefacesection{Plagiarism statement}
\bigskip \noindent I declare that this thesis is my
own work, except where acknowledged, and has not been submitted for
academic credit elsewhere. 

\bigskip \noindent I acknowledge that the assessor of this
thesis may, for the purpose of assessing it:
\begin{itemize}
\item Reproduce it and provide a copy to another member of the University; and/or,
\item Communicate a copy of it to a plagiarism checking service (which may then retain a copy of it on its database for the purpose of future plagiarism checking).
\end{itemize}

\bigskip \noindent I certify that I have read and understood the University Rules in
respect of Student Academic Misconduct, and am aware of any potential plagiarism penalties which may 
apply.\vspace{24pt}

\bigskip \noindent By signing 
this declaration I am
agreeing to the statements and conditions above.

\bigskip \noindent
Signed: \rule{7cm}{0.25pt} \hfill Date: \rule{4cm}{0.25pt} 


% Acknowledgements
% \pagebreak
% \prefacesection{Acknowledgements}
% \bigskip \noindent Body text


% Abstract
% \pagebreak
% \prefacesection{Abstract}
% \bigskip \noindent Body text

\afterpreface


%===============================================================================================================================%
%                                                          INTRODUTION                                                          %
%===============================================================================================================================%

\pagebreak
\chapter*{Introduction}\label{s-intro}
\addcontentsline{toc}{chapter}{Introduction}

The goal of this paper is to develop methods for factoring large integers using the theory of binary quadratic forms.
The paper is written with the intention of providing motivated explanations to the reader,
as the existing resources often give proofs of theorems without explaining how such a result would be formulated without prior knowledge of its truth.
This poses a particular challenge in \Cref{s-comp}, 
as I'm yet to find a paper that approaches composition in the style of a derivation instead of stating it as a result and then proving that it works.

\vspace{2mm}
\noindent
For the chapter draft, 
\begin{itemize}
    \item A completed draft of \Cref{s-forms} is included, where some of the basic definitions are established.
    Here, we establish that forms of a fixed discriminant can be partitioned into a finite set of equivalence classes.

    \item A bare-bones version of \Cref{s-comp} is included.
    Here, important results are given, and the direction of the proofs are in place, but many of the details have not yet been formalised.
    Notes have been left in \red{red text} to indicate work that is needed.

    \item The beginning of \Cref{s-factor} is also included.
    Here, the usefulness of binary quadratic forms for factoring is realised.
    The key result that the previous chapters have built towards is provided in \Cref{s-factor-amb};
    ``ambiguous forms'' immediately provide a factorisation of the discriminant.
\end{itemize}

The tasks for the remainder of the project are to finish formalising the derivations in \Cref{s-comp},
and to develop methods for finding ambiguous forms efficiently.

So far, the paper has focuseed on definite forms (negative discriminant), 
as they have a simpler structure and are sufficient to develop factoring algorithms.
However, at some point we will likely want to use definite forms (positive discriminant), as they give rise to different algorithms \cite[\S\,8.7]{cohen}.
This will require adding to the content of \Cref{s-forms}, and possibly adjusting the arguments in \Cref{s-comp}.

Beyond this, if time allows, 
we could move to an investigation of factoring algorithms which leverage larger number fields than the quadratic fields which underpin binary quadratic forms.

Alternatively, we could move to a discussion of cryptography \cite[Ch.\,12]{buchmann};
if finding ambiguous forms leads immediately to factoring the discriminant, then mapping the class group is at least as computationally complex as factoring.
Since existing cryptographic methods depend on the assumption that factoring is difficult,
cryptographic methods which depend on the structure of the class group should, in principal, be secure.
Developing a novel encryption method may be an interesting way to continue the narrative of the paper,
but whether this will be an appropriate use of time remains to be seen.


%===============================================================================================================================%
%                                                           CHAPTER 1                                                           %
%===============================================================================================================================%

\pagebreak
\chapter{Binary Quadratic Forms}\label{s-forms}

\hspace{-3.5mm}
\textit{
    In this chapter, we establish standard definitions and results on binary quadratic forms,
    as can be found in \cite[Ch.\,1,\,2,\,5]{buchmann}, \cite[Ch.\,1-2]{buell}, \cite[\S\,2]{cox}.
}

\vspace{2mm}
\noindent
A \textbf{binary quadratic form} (or simply a \textbf{form}) is a function of the form
$$f(x, y) = ax^2 + bxy + cy^2.$$
It is called ``binary'' because it is a function of two variables, and ``quadratic'' because the degree of each term is 2.
For our purposes, we are usually working over the integers. 
Since $x$ and $y$ are variables, the behaviour of a given form is wholly dependent on the values of $a$, $b$ and $c$,
which motivates a shorthand notation of the same form; $(a, b, c)$.
We will impose that $a, c \neq 0$, or else $x$ or $y$ would be reduced to a linear factor.
Allowing this would serve no purpose beyond comlicating proofs with the added work of edge case considerations.

A form is called \textbf{primitive} if $\text{gcd}(a, b, c) = 1$. 
Primitive forms are of primary concern, as he behaviour of a nonprimitive form can be more simply understood by removing the common factor to work with simpler coefficients.


\section{Representation}\label{s-forms-rep}

A form $f(x, y)$ is said to \textbf{represent} a number $n$ if there exists an integer solution
$$f(x_0, y_0) = n \rnote{x_0, y_0 \in \Z.}$$
For example, consider the form $(2, 1, 5)$, or $f(x, y) = 2x^2 + xy + 5y^2$. 
This form represents the number 2, since $f(1, 0) = 2$.
The form does not represent the number 1, though it is not immediately clear how to prove a result of this kind.

The questions which historically motivated the study of forms were as follows \cite[Ch.\,1]{buell};
which numbers can be represented by a given form and, conversely, which forms can represent a given number?
These questions will not be our primary concern, but they will help to motivate the development of the theory that we need.


\pagebreak
\section{Discriminant}\label{s-forms-disc}

The form $x^2 + y^2$ clearly exclusively represents non-negative integers. 
The same property holds for the previous example, $2x^2 + xy + 5y^2$, though this is harder to see.
On the other hand, the form $x^2 - y^2$ clearly represents both positive and negative integers.
A reasonable first step to understanding a form's representations is to classify whether they can represent positive, negative, or both positive and negative integers.
To do this, let us consider a related continuous, real-valued function, 
$$g(t) = at^2 + bt + c.$$
Notice that
\begin{align*}
    g\left(\frac{x}{y}\right) &= a\left(\frac{x}{y}\right)^2 + b\left(\frac{x}{y}\right) + c \\
    &= \left(ax^2 + bxy + cy^2\right)\frac{1}{y^2} \\
    &= f(x, y) \frac{1}{y^2}.
\end{align*}
Hence we see that $f(x, y)$ and $g(\tfrac{x}{y})$ agree up to sign, with the exception of when $y=0$.
Now, the zeroes of $g(t)$ are given by
$$t = \frac{-b \pm \sqrt{b^2 - 4ac}}{2a}.$$
If $b^2 -4ac$ is negative, then $g$ has no real zeroes, meaning that $g$ must be either positive everywhere or negative everywhere (except the origin).
Since $f(x, y)$ agrees up with $g(\tfrac{x}{y})$ up to sign, then $f$ must also have a fixed sign for its entire range of outputs - except possibly when $y=0$.
However, note that $f$ can be extended to a continuous function over $\R^2$ and that $f(x, 0) = ax^2$, which can only be zero when $x$ is zero (since $a \neq 0$).
Combining these facts, we can argue that the line $y=0$ must have the same sign as the rest of function, with the exception of the trivial zero.

We call forms satisfying $b^2 -4ac < 0$ \textbf{positive definite} or \textbf{negative definite}, depending on whether they output positive or negative values.
For forms where $b^2 - 4ac > 0$, $g$ has two distinct real zeroes, so $g$ (and hence $f$) will be positive for some inputs and negative for others.
Such forms are called \textbf{indefinite}.
Similar reasoning shows that forms with $b^2 -4ac = 0$, called \textbf{degenerate forms}, are strictly non-negative or non-positive, but such forms are of little intereest.

This motivates defining the \textbf{discriminant} of a form $ax^2 + bxy + cy^2$ to be 
$$\Delta = b^2 - 4ac.$$
We will come to see that this classification is far more powerful than simply being able to classify forms as definite or indefinite.

\pagebreak
\begin{lemma}
    The discriminant of a form, $\Delta$, is equivalent to 1 or 0 modulo 4. 
\end{lemma}

\begin{proof}
    The definition of the discriminant gives
    $$\Delta \equiv b^2 \pmod{4}$$
    and the only possible values for a perfect square modulo 4 are 1 and 0.
\end{proof}

\begin{lemma}\label{lemma-discriminant-under-multiple}
    If the discriminant of the form $(a, b, c)$ is $\Delta$, then the discriminant of the form $(ak, bk, ck)$ is $\Delta k^2$.
\end{lemma}

\begin{proof}
    The discriminant of $(ak, bk, ck)$ is given By
    $$(bk)^2 - 4(ak)(ck) = b^2k^2 - 4ack^2 = \Delta k^2.$$
\end{proof}


\section{Equivalence}\label{s-forms-eq}
Given a form $f(x, y) = ax^2 + bxy + cy^2$, one can generate new forms which represent the same set of integers by changing the input space;
take some invertible function $\phi : \Z^2 \rightarrow \Z^2$, and let $g(x, y) = f \circ \phi(x, y)$.
Then for any $x_0, y_0 \in \Z$
$$f(x_0, y_0) = n \implies g(\phi(x_0, y_0)) = n$$
$$g(x_0, y_0) = m \implies f(\phi^{-1}(x_0, y_0)) = m.$$
That is, all integers represented by $f$ are also represented by $g$, and vice versa. 
However, we must also ensure that our choice of $\phi$ gives $g$ to be a form.
Consider the choice $\phi(x, y) = (x+1, y)$.
This is an invertible function on $\Z^2$, but when it is composed with our usual choice of $f$ we have
\begin{align*}
    g(x, y) &= f \circ \phi(x, y) \\
    &= f(x+1, y) \\
    &= a(x+1)^2 + b(x+1)y + cy^2 \\
    &= ax^2 + 2ax + a + bxy + by + cy^2.
\end{align*}
This is no longer a quadratic form, so this choice of $\phi$ should not be allowed. 
Instead, consider a function of the form $\phi(x, y) = (px+qy, rx+sy)$.
This gives
\begin{align*}
    &\hspace{-5mm} g(x, y) = f \circ \phi(x, y) \\
    &= f(px+qy, rx+sy) \\
    &= a(px+qy)^2 + b(px+qy)(rx+sy) + c(rx+sy)^2 \\
    &= \left[ap^2 + bpr + cr^2 \right] x^2
        + \left[ 2apq + b(ps + qr) + 2crs \right] xy
        + \left[ aq^2 + bqs + cs^2 \right] y^2
\end{align*}
which gives $g$ to be a binary quadratic form, as required. 
We now require that $\phi$ is invertible, so that this is a symmetric relationship.
Consider that $\phi$ can be understood as a matrix multiplication
$$\phi \left( \bmat{x \\ y} \right) = \bmat{p & q \\ r & s} \bmat{x \\ y}.$$
So it is clear to see that the valid choices of $\phi$ are the invertible elements of $M_2(\Z)$. 
These are the elements of $\GL_2(\Z)$, also known as the \textbf{unimodular} matrices (of size 2).
The determinant of such a matrix will always be $\pm 1$ 
(since the determinant is a result of sums and products of the integer coefficients, 
and the product of determinants of a matrix and its inverse is 1).

We define forms to be \textbf{improperly equivalent} if there is a $\GL_2(\Z)$ transformation between them.
It is easy to see that this is an equivalence realtion.
If we restrict the discriminant of the transformation to be strictly positive 1, we split the equivalence classes further-
we call such forms \textbf{properly equivalent}, which is also an equivalence relation.

Let us consider the effect of $\SL_2(\Z)$ transformations on forms;
for $f \circ S = (c, -b, a)$, the discriminant is clearly unchanged.
For $f \circ T = (a, 2a + b, a + b + c)$, we have
\begin{align*}
    \Delta &= (2a + b)^2 - 4a(a + b + c) \\
    &= 4a^2 - 4ab + b^2 - 4a^2 + 4ab - 4ac \\
    &= b^2 - 4ac.
\end{align*}
So the discriminant is also unchanged under the action of $T$.
Hence we have

\vspace{2mm}
\begin{lemma}
    Properly equivalent forms have equal discriminants.
\end{lemma}

\vspace{2mm}
\begin{remark}
    Properly/improperly equivalent forms represent exactly the same set of integers.
\end{remark}

\begin{proof}
    This is by design of our formulation of equivalence.
    Let us now prove it formally.
    We will prove the case of improper equivalence, as any forms which are properly equivalent are also improperly equivalent.

    Suppose $f = (a, b, c)$ and $g = (a', b', c')$ are improperly equivalent forms.
    That is, we have some $\phi \in \GL_2(\Z)$ such that
    $$f \circ \phi(x, y) = g(x, y).$$
    If $g$ represents $n$, then there exist integers $x_0, y_0 \in \Z$ such that
    \begin{align*}
        g(x_0, y_0) &= n \\
        \implies f \circ \phi(x_0, y_0) &= n.
    \end{align*}
    But $\phi$ is a function from $\Z^2 \rightarrow \Z^2$, so we see that $f$ also represents $n$.
    The converse argument, $f$ represents $n$ $\implies$ $g$ represesnts $n$, is identical.
    Hence, $f$ represents $n$ if and only if $g$ represents $n$.
\end{proof}


\pagebreak
\section{Reduction}\label{s-forms-red}
Now that we have an equivalence relation, we would like to be able to find a suitable representative element of a given class.
This will not only prove useful for simplifying computations by choosing a representative with small values of $a$, $b$, and $c$,
but will also give us a method for checking whether two forms belong to the same class by evaluating and comparing their canonical representatives.

A basis for $\SL_2(\Z)$ is given by
$$S = \pmat{0 & -1 \\ 1 & 0}, \hspace{5mm} T = \pmat{1 & 1 \\ 0 & 1}.$$
Note that
$$S^4 = I, \hspace{5mm} T^k = \pmat{1 & k \\ 0 & 1}.$$
So we see that any properly equivalent forms must be related by some series of applications of $S$ and $T^k$.
Let us observe the results of these transformations on a generic form $f(x, y) = (a, b, c)$
\begin{align*}
    f \circ S (x, y) &= f \circ \pmat{0 & -1 \\ 1 & 0} \pmat{x \\ y} \\
    &= f(-y, x) \\
    &= a(-y)^2 + b(-y)x + cx^2 \\
    &= ay^2 - bxy + cx^2
\end{align*}
\begin{align*}
    f \circ T^k (x, y) &= f \circ \pmat{1 & k \\ 0 & 1} \pmat{x \\ y} \\
    &= f(x + ky, y) \\
    &= a(x + ky)^2 + b(x + ky)y + cy^2 \\
    &= a(x^2 + 2kxy + k^2y^2) + b(xy + ky^2) + cy^2 \\
    &= ax^2 + (2ak + b)xy + (ak^2 + bk + c)y^2.
\end{align*}
So we have
\begin{gather*}
    S : (a, b, c) \mapsto (c, -b, a) \\
    T^k : (a, b, c) \mapsto (a, 2ak + b, ak^2 + bk + c).
\end{gather*}
Let us use these for constructing a redution algorithm, which turns any form into a canonical class representative.
For now, we will only consider positive definite forms, as the result is simpler.
Note that in this case, $a$ and $c$ must both be positive, as their signs must be equal for $\Delta = b^2 - 4ac < 0$,
and they must be positive to represent positive integers.

A natural way to use $S$ and $T$ to reduce a form is to attempt to minimise $|b|$.
Whenever $|b| > a$, we can use $T$ to reduce $b$ to be in the range $[-a, a]$.
Then, if $c < a$, we can use $S$ to swap $a$ and $c$, then return to step 1 to reduce $b$ further.
Eventually, we reach a form with $|b| \leq a \leq c$, at which point the algorithm terminates.

In the case that $a = |b|$ or $a = c$, a form has two representations which satisfy $a \leq b \leq c$ --- one with positive $b$, and one with negative $b$.
This can be fixed by taking $b$ to be positive in these cases.
We adjust the reduction algorithm such that step 1 moves $b$ into the interval $(-a, a]$ (removing the choice of sign in the case that $|b| = a$),
and applying $S$ not only if $a > c$, but if $a = c$ and $b$ is negative.
The full reduction algorithm is as follows

\begin{algorithm}\label{alg-reduction}
    \caption{Reduction of positive definite forms}
    \DontPrintSemicolon
    \KwData{$a, b, c \in \Z$}
    \While {$b \notin (-a, a]$}{
        $k \gets \lceil (b - a) / 2a \rceil$\;
        $c \gets k^2 a + k b + c$\;
        $b \gets 2ka + b$\;
        \If {$a > c$ or ($a = c$ and $b < 0$)}{
            $a \leftrightarrow c$ \hspace{2mm} (swap values)\;
            $b \gets -b$\;
        }
    }
\end{algorithm}
\noindent
So we say that a positive definite form is \textbf{reduced} if it satisfies
\begin{itemize}
    \item $|b| \leq a \leq c$
    \item $b \geq 0$ if $a = |b|$ or $a = c$.
\end{itemize}

\vspace{2mm}
\begin{theorem}
    \Cref{alg-reduction} terminates in a finite number of steps.
\end{theorem}

\begin{proof}
    On each iteration of the while loop, the value of $b$ decreases.
    Since $b$ is initially a finite integer, and only takes on integer values, the initial value of $b$ is an upper bound on how many loops the algorithm can perform before terminating.
\end{proof}

\vspace{2mm}
\begin{theorem}
    Every positive definite form is properly equivalent to exactly one reduced form. 
\end{theorem}

\begin{proof}
    The fact that any form is equivalent to at least one reduced form follows immediately from the existence of the reduction algorithm.
    To prove uniqueness, we adapt the proof of Cohen \cite[Prop.\,5.3.3]{cohen}.
    Suppose $f = (a, b, c)$ is a reduced form, and that $f' = (a', b', c')$ is another form, such that these forms are properly equivalent.
    That is, there exists a transformation 
    $$\phi = \pmat{p & q \\ r & s} \in \SL_2(\Z).$$
    such that $f \circ \phi = f'$.
    We will show that if $f'$, then it must be equal to $f$.

    First, we claim that $a$ is the minimal coefficient of $x^2$ for all forms in its equivalence class --- that is, $a \leq a'$.
    In \Cref{s-forms-eq}, we showed that the result of $\phi$ acting on $f$ is
    $$\left[ap^2 + bpr + cr^2 \right] x^2
        + \left[ 2apq + b(ps + qr) + crs \right] xy
        + \left[ aq^2 + bqs + cs^2 \right] y^2.$$
    That is, $a' = ap^2 + bpr + cr^2$.
    Using the fact that $f$ is reduced, we have $c \geq a$ and $b > -a$.
    Hence, we can bound $a'$;
    $$a' = ap^2 + bpr + cr^2 \geq ap^2 - a|pr| + ar^2$$
    Since $b$ is strictly greater than $-a$, equality can only hold when $pr = 0$.
    Simplifying this result yields
    \begin{align*}
        a' &\geq a \left[ |p|^2 - |pr| + |r|^2 \right] \\
        &= a \left[ (|p| - |r|)^2 + |pr| \right].
    \end{align*}
    It's impossible for both $pr$ and $|p| - |r|$ to be zero, as this would require $p = r = 0$,
    which would cause the determinant of $\phi$ to be zero.
    But the determinant of $\phi$ is known to be 1, so, since $p, r \in \Z$, we have that $a'$ is at least $a$.

    Now, consider the case where $f'$ is reduced.
    Then by the same argument as above, we have that $a' \geq a$ --- together these results imply $a = a'$.
    That is, properly equivalent reduced forms must have the same leading coefficient.
    
    Now we must show that if $\phi$ fixes $a$, then the action of $\phi$ on $f$ is trivial.
    Recall that in order for the equality $a' = a$ to hold, we must have $pr = 0$.
    This gives
    $$a' = a \left[ p^2 + r^2 \right]$$
    So we also require $p^2 + r^2 = 1$.
    This leaves us with two possible cases
    \begin{gather*}
        \textbf{(1)} \hspace{2mm} p = \pm 1, r = 0 \\
        \textbf{(2)} \hspace{2mm} p = 0, r = \pm 1.
    \end{gather*}
    Consider the first case, letting $p = \delta \in \{1, -1\}$.
    Since $\phi \in \SL_2(\Z)$, we have
    $$1 = \text{det}(\phi) = ps - qr = \delta s$$
    which gives $s = \delta$.
    Hence we have
    $$\phi = \pmat{\delta & q \\ 0 & \delta} = \pmat{\delta & 0 \\ 0 & \delta} T^q.$$
    Note the action of the matrix
    $$\pmat{-1 & 0 \\ 0 & -1} : (a, b, c) \mapsto (a, b, c)$$
    so $\phi$ has the same action as $T^q$. 
    Since this fixes $a$ and shifts $b$ by $2qa$, any choice for $q$ other than 0 (which gives the identity)
    shifts $b$ outside of the range $(-a, a]$, giving a form which is no longer reduced.
    That is, $f' = f \circ \phi = f$.

    In the second case, let $r = \delta \in \{1, -1\}$. 
    The equality $a' = a$ depends on $cr^2 \geq ar^2$ acheiving equality. 
    Since $r$ is nonzero and $c \geq a$, this is only acheived if $c = a$.
    Again taking the determinant of $\phi$, we have
    $$1 = \text{det}(\phi) = ps - qr = - \delta q$$
    which gives $q = - \delta$.
    Hence we have
    \begin{align*}
        \phi &= \pmat{0 & - \delta \\ \delta & s} \\ 
        &= \pmat{0 & -1 \\ 1 & 0} \pmat{\delta & 0 \\ 0 & \delta} \pmat{1 & s \\ 0 & 1} \\
        &= S \cdot \pmat{\delta & 0 \\ 0 & \delta} \cdot T^s.
    \end{align*}
    As before, we can ignore the action of the negative identity matrix.
    Hence we see that the action of $\phi$ is given by
    $$ST^s : (a, b, a) \mapsto (as^2 + bs + a, - 2as - b, a).$$
    Recall that the leading term must be equal to $a$ in order for this form to be reduced.
    Then we have
    \begin{gather*}
        as^2 + bs + a = a \\
        \implies 0 = as^2 + bs = s(as + b)
    \end{gather*}
    So either $s = 0$ or $as + b = 0$. 
    In the former case, the action of $\phi = S$ flips the sign of the second coefficient.
    If $b \neq 0$, this yields a form that is no longer reduced, and if $b = 0$, it acts trivially.
    
    In the latter case, since $(a, b, a)$ is reduced, we have $a \geq b$ and $|b| \geq 0$.
    To satisfy $as + b = 0$ for $a, b \in \Z$ and $s = \pm 1$, we must therefore have $a = b$, $s = -1$.
    The action of $\phi$ is therefore
    $$\phi : (a, a, a) \mapsto (a, a, a).$$
    Hence we see that the only reduced form which is properly equivalent to $(a, b, c)$ is itself.
\end{proof}

It seems natural to ask, how many reduced forms of a given discriminant exist?
That is, how many equivalence classes of a fixed discriminant are there?
For a reduced form, we have a bound on $b$ given by $|b| \leq a \leq c$.
Does this truly allow us to bound $b$?
We need to use the fact that the choices for $a$ and $c$ are constrained by the discriminant:
\begin{align*}
    b^2 - 4ac &= \Delta \\
    b^2 - 4 \cdot |b| \cdot |b| &\geq \Delta \\
    -3b^2 &\geq \Delta \\
    b^2 &\leq -\tfrac{\Delta}{3} \\
    |b| &\leq \sqrt{-\tfrac{\Delta}{3}}
\end{align*}
so there are only finitely many choices for $b$ in a reduced positive definite form.
If we fix a value of $b$, then can we bound the number of choices for $a$ and $c$?
Consider
\begin{align*}
    b^2 - 4ac &= \Delta \\
    -4ac &= \Delta - b^2 \\
    ac &= \frac{b^2 - \Delta}{4}
\end{align*}
so we have finitely many choices for $a$ and $c$ --- each corresponding to a factorisation of $\tfrac{b^2 - \Delta}{4}$.
Hence, we see that there are only finitely many choices for $a$, $b$, and $c$ when constructing a reduced positive definite form of a fixed discriminant.
We have proved

\vspace{2mm}
\begin{theorem}
    The set of equivalence classes of positive definite forms of fixed discriminant $\Delta < 0$ is finite.
\end{theorem}


%===============================================================================================================================%
%                                                           CHAPTER 2                                                           %
%===============================================================================================================================%

\chapter{Composition}\label{s-comp}
Now that we have the means to properly understand equivalence classes of forms, we seek to equip them with a binary operation to form a group.
We acheive this by showing a bijection between classes of forms and classes of ideals, 
under which the ideal multiplication structure gives a multiplicative structure to the corresponding classes of forms.

\section{Historical Background}
\emph{In this section, we draw from \cite[\S2D-\S3A]{cox}.}

\vspace{2mm}
\noindent
Legendre was the first to develop a notion of composition, inspired by Lagrange's use of the identity
$$(2x^2 + 2xy + 3y^2)(2z^2 + 2zw + 3w^2) = (2xz + xw + yz + 3yw)^2 + 5(xw - yz)^2.$$
We see that for any two numbers which can be represented by the form $(2, 2, 3)$, their product can be represented by the form $(1, 0, 5)$.
That is, if we have representations of $n$ and $m$
\begin{gather*}
    2x^2 + 2xy + 3y^2 = n \\
    2z^2 + 2zw + 3w^2 = m
\end{gather*}
then the form $(1, 0, 5)$ represents $nm$ via
$$(2xz + xw + yz + 3yw)^2 + 5(xw - yz)^2 = nm.$$
Legendre sought to generalise results of this kind.
He defined a form $F$ to be the \textbf{composition} of $f$ and $g$ if there exist bilinear forms $B_1$ and $B_2$ such that
\begin{gather*}
    f(x, y)g(z, w) = F(B_1(x, y; z, w), B_2(x, y; z, w)) \\
    B_i = a_i xz + b_i xw + c_i yz + d_i yw
\end{gather*}
Legendre was able to show that any two forms of the same discriminant can always be composed in this way, but his formulation was multi-valued, giving several options for the resulting form $F$.
This was finally rectified by Gauss in his \emph{Disquisitiones Arithmaticae} \cite{gauss}, where he refined Legendre's formulation to turn composition of forms into a well-defined function.
This notion of composition gave a group structure to the set of proper equivalence classes of forms of a fixed discriminant.
This group is called the \textbf{class group}.

Due to the complexity of Gauss' methods, Dirichlet presented a simpler but less general approach to composition.
We will acquire a similar formulation of composition in this chapter.
However, as the theory of ideals developed, it became apparent that there exists a bijection between certain ideals of quadratic number fields and the class group of quadratic forms \cite[Ch.\,1]{buell}.
This provides us a new and perhaps more intuitive way to derive composition of forms;
we first outline this bijection in \Cref{s-comp-ideals}, then we can leverage the multiplicative structure of ideals to define an analogous operation on forms in \Cref{s-comp-formula}.
We will then show that this operation precisely fits the desired definition of composition given by Legendre.

\section{Algebraic Number Theory Preliminaries}\label{s-comp-algebraic}
For a finite field extension of the rational numbers, $K/\Q$, an \textbf{order} in $K$ is some $\order \subseteq K$ if \cite[Def.\,4.6.1]{cohen}
\begin{itemize}
    \item $\order$ is a subring of $K$
    \item $\order$ is a finitely generated $\Z$-module
    \item the $\Q$-span of $\order$ is $K$ (ie. $\text{rank}(\order) = \deg(K:\Q)$)
\end{itemize}
Here, we are working with the definition of a ring that requires the existence of a multiplicative identity.

Now, for $\Delta \equiv 0, 1 \pmod{4}$,
Let $K_\Delta = \Q(\sqrt\Delta)$, and define the \textbf{order of discriminant} $\mathbf{\Delta}$ to be 
$$\order_\Delta = \Z + \frac{\Delta + \sqrt\Delta}{2}\Z$$

\begin{lemma}
    $\order_\Delta$ is an order in $K_\Delta$.
\end{lemma}

\begin{proof}
    The latter two conditions of being an order are met by definition, so we need only show that $\order_\Delta$ is a subring of $K_\Delta$. 
    It is also clear that $\order_\Delta \subseteq K_\Delta$, so if we show that $\order_\Delta$ is closed under ring operations, then we can conclude that it is an order in $K_\Delta$.
    Closure under addition is obvious - as for closure under multiplication, let us show closure holds when multiplying the generators of $\order_\Delta$ with themselves, from which closure of all elements under multiplication can be inferred.
    The generators of $\order_\Delta$ are 1 and $\frac{\Delta + \sqrt\Delta}{2}$.
    Closure under multiplication by 1 is obvious, so we need only check
    \begin{align*}
        \frac{\Delta + \sqrt\Delta}{2} \cdot \frac{\Delta + \sqrt\Delta}{2} &= \frac{\Delta^2 + 2\Delta\sqrt\Delta + \Delta}{4} \\
        &= \frac{\Delta - \Delta^2}{4} + \Delta\frac{\Delta + \sqrt\Delta}{2} \\
        &= \frac{\Delta (1 - \Delta)}{4} + \Delta\frac{\Delta + \sqrt\Delta}{2}
    \end{align*}
    and since $\Delta \equiv 0, 1 \pmod{4}$, the first term is an integer, so this result is in $\order_\Delta$.
\end{proof}

\begin{lemma}\label{lemma-ideal-lattice}
    Any nonzero ideal $I \trianglelefteq \order_\Delta$ can be written in the form
    $$I = l\Z + \left( m + n\frac{\Delta + \sqrt\Delta}{2} \right)\Z$$
    where $l, m, n \in \Z$ and $l, n > 0$ are uniquely determined.
\end{lemma}

\pagebreak
\begin{proof}
    We will first show that $I$ contains at least one integer element, and at least one element with a nonzero $\sqrt\Delta$ component.
    Then, we choose $l, m$ and $n$ according to certain minimum conditions, and argue that their $\Z$-span is all of $I$, or else a contradiction would arise.
    
    Consider a nonzero element $x + \frac{\Delta + \sqrt\Delta}{2}y \in I$. 
    The conjugate of this element is
    $$x + \frac{\Delta - \sqrt\Delta}{2}y = x + y\Delta - \frac{\Delta + \sqrt\Delta}{2}y \in \order_\Delta.$$
    Since $I$ is an ideal of $\order_\Delta$, multiplication by an element of $\order_\Delta$ must give a result which is still in $I$.
    In particular, multiplying our element by its conjugate gives
    \begin{align*}
        I &\ni \left( x + \frac{\Delta + \sqrt\Delta}{2}y \right) \left( x + \frac{\Delta - \sqrt\Delta}{2}y \right) \\
        &= \left( x + \frac{\Delta}{2}y \right)^2 -  \left( \frac{\sqrt\Delta}{2}y \right)^2 \\
        &= x + \Delta xy + \frac{\Delta^2 y^2}{4} - \frac{\Delta y^2}{4} \\
        &= x + \Delta xy + \frac{\Delta(\Delta - 1)y^2}{4}.
    \end{align*}
    This is an integer, so there is at least one integer in $I$.
    Multiplying this by $\frac{\Delta + \sqrt\Delta}{2} \in \order_\Delta$, we see that there is at least one element in $I$ with a nonzero $\frac{\Delta + \sqrt\Delta}{2}$ component.
    
    Now, choose $l$ to be the minimal positive integer in $I$.
    Every integer $z \in I \cap \Z$ must ne a multiple of $l$, or else we would have $l > \gcd(l, z)$ --- a contradiction with the choice of $l$.
    We can therefore see that the choice of $l$ is unique (up to sign, but we take $l$ to be positive).

    Likewise, choose $m + n\frac{\Delta + \sqrt\Delta}{2} \in I$ such that $n > 0$ is minimum.
    For any $x + \frac{\Delta + \sqrt\Delta}{2}y \in I$ we must also have $n \divides y$, or else a linear combination of elements would yield a result whose coefficient of $\frac{\Delta + \sqrt\Delta}{2}$ would be $\gcd(y, n)$, contradicting the choice of $n$.
    This uniquely determines $n$ (again, up to sign, but we take $n > 0$).

    Finally, we will show that our general element can be expressed as an element of $l\Z + \left( m + n\frac{\Delta + \sqrt\Delta}{2} \right)\Z$.
    Letting $y = y_0 n$, we can write
    \begin{align*}
        x + \frac{\Delta + \sqrt\Delta}{2}y &= x - y_0 m + y_0 m + \frac{\Delta + \sqrt\Delta}{2}y_0 n \\
        &= x - y_0 m + \left( m + n\frac{\Delta + \sqrt\Delta}{2} \right) y_0.
    \end{align*}
    Now, $(m + n\frac{\Delta + \sqrt\Delta}{2}) y_0 \in I$, so by closure under addition, $x - y_0 m$ must also be in $I$.
    This is an integer, so it can be written as a multiple of $l$.
    Hence $x + \frac{\Delta + \sqrt\Delta}{2}y \in l\Z + (m + n\frac{\Delta + \sqrt\Delta}{2})\Z$, so $I = l\Z + (m + n\frac{\Delta + \sqrt\Delta}{2}\Z)$.
\end{proof}

\pagebreak
\begin{remark}
    We could also choose a canonical value of $m \in [0, l)$, by performing a change of basis, adding a multiple of $l$ to this second generator.
    However, this would only create further work in later proofs when we'd have to ensure that $m$ is in this interval --- so we don't enforce this condition.
\end{remark}

\vspace{2mm}
\begin{lemma}\label{lemma-lattice-divisibility-conditions}
    The $\Z$-module
    $I = l\Z + \left( m + n\frac{\Delta + \sqrt\Delta}{2} \right)\Z$
    is an ideal of $\order_\Delta$ if and only if the following conditions are met:
    \begin{enumerate}
        \item $n \divides m$
        \item $n \divides l$
        \item $l \divides \left( \frac{m^2}{n} + m\Delta + n \frac{\Delta - \Delta^2}{4} \right).$
    \end{enumerate}
\end{lemma}

\begin{proof}
    First, suppose that $I = l\Z + \left( m + n\frac{\Delta + \sqrt\Delta}{2} \right)\Z$ is an ideal of $\order_\Delta$.
    This means that $I$ is closed under multiplication by elements of $\order_\Delta$.
    In particular,
    $$l \cdot \frac{\Delta + \sqrt\Delta}{2}, \left( m + n\frac{\Delta + \sqrt\Delta}{2} \right) \cdot \frac{\Delta + \sqrt\Delta}{2} \in I.$$
    Firstly,
    \begin{align*}
        l \cdot \frac{\Delta + \sqrt\Delta}{2} &= \left( n\frac{\Delta + \sqrt\Delta}{2} \right) \frac{l}{n} \\
        &= - \frac{m}{n}l + \frac{l}{n}\left( m + n\frac{\Delta + \sqrt\Delta}{2} \right) \in I.
    \end{align*}
    Since $I$ is a free $\Q$-module, representation of elements as linear combinations of basis elements are unique.
    Hence, for this to be in $I$ we require that the coefficients, $-\frac{m}{n}$ and $\frac{l}{n}$, be integers.
    Hence, we require $n \divides m$ and $n \divides l$. Secondly, we have
    \begin{align*}
        &\hspace{-5mm} \left( m + n\frac{\Delta + \sqrt\Delta}{2} \right) \cdot \frac{\Delta + \sqrt\Delta}{2} \\
        &= m \frac{\Delta + \sqrt\Delta}{2} + n \frac{\Delta^2 + 2\Delta\sqrt\Delta + \Delta}{4} \\
        &= m \frac{\Delta + \sqrt\Delta}{2} + n \Delta \frac{\sqrt\Delta}{2} + n \frac{\Delta^2 + \Delta}{4} \\
        &= (m + n\Delta)\frac{\Delta + \sqrt\Delta}{2} - n \frac{\Delta^2}{2} + n \frac{\Delta^2 + \Delta}{4} \\
        &= n\left( \frac{m + n\Delta}{n} \right)\frac{\Delta + \sqrt\Delta}{2} + n \frac{\Delta - \Delta^2}{4} \\
        &= \left( \frac{m + n\Delta}{n} \right) \left(m + n\frac{\Delta + \sqrt\Delta}{2} \right) - m(\frac{m + n\Delta}{n}) + n \frac{\Delta - \Delta^2}{4} \\
        &= \left( \frac{m}{n} + \Delta \right) \left(m + n\frac{\Delta + \sqrt\Delta}{2} \right) - \frac{1}{l}\left( \frac{m^2}{n} + m\Delta + n \frac{\Delta - \Delta^2}{4} \right)l.
    \end{align*}
    Again, for this to be in $I$, we require the coefficients to be integers.
    Hence we have the added condition $l \divides \frac{m^2}{n} + m\Delta + n \frac{\Delta - \Delta^2}{4}$.

    Conversely, let $I = l\Z + \left( m + n\frac{\Delta + \sqrt\Delta}{2} \right)\Z$ satisfy the given divisibility conditions.
    $I$ is a commutative $\Z$-module, so it is closed under addition and under multiplication by integers.
    By the calculations above, we see that the given divisibility conditions are sufficient for $l$ and $m + n\frac{\Delta + \sqrt\Delta}{2}$ to remain in $I$ when multiplied by $\frac{\Delta + \sqrt\Delta}{2}$.
    Hence, a generic product of elements in $I$ and $\order_\Delta$ is
    \begin{align*}
        &\hspace{-5mm} \left[ lx + \left( m + n\frac{\Delta + \sqrt\Delta}{2} \right)y \right] \cdot \left( x' + \frac{\Delta + \sqrt\Delta}{2} y' \right) \rnote{x, x', y, y' \in \Z} \\
        &= lxx' + \left( m + n\frac{\Delta + \sqrt\Delta}{2} \right) yx' + l \frac{\Delta + \sqrt\Delta}{2} xy' + \left( m + n\frac{\Delta + \sqrt\Delta}{2} \right) \frac{\Delta + \sqrt\Delta}{2} yy'.
    \end{align*}
    Each term is in $I$, so their sum is in $I$. 
    Hence, $I$ is an ideal of $\order_\Delta$.
\end{proof}

We also need to define the norm of an ideal.
For an ideal $I$ of a ring $R$, the \textbf{norm} is defined as
$$N(I) = \left| R / I \right|.$$

\begin{lemma}\label{lemma-ideal-norm}
    The norm of $I = l\Z + \left( m + n\frac{\Delta + \sqrt\Delta}{2} \right) \Z \trianglelefteq \order_\Delta$ is given by
    $$ N(I) = l \cdot n.$$
\end{lemma}

\begin{proof}
    Consider the elements of the quotient $\order_\Delta / I$.
    Claim that these are given by
    $$\order_\Delta / I = \left\{ x + \frac{\Delta + \sqrt\Delta}{2}y \mathrel{:} 0 \leq x < l, 0 \leq y < n \right\}.$$
    To see that this is the correct formulation, consider a generic element $\order_\Delta$.
    By subtracting multiples of $m + n\frac{\Delta + \sqrt\Delta}{2}$ and then multiples of $l$, we can uniquely determine representative for this element's class such that $0 \leq y < n$, and $0 \leq x < l$.
\end{proof}

Finally, we will define the norm of an element of a number field.
This can be defined more generally for larger fields \todo{cite definition}, but for our purposes, we will only handle the quadratic case.
For any quadratic number field $K$ (that is, $K = \Q(\sqrt{\Delta})$ where $\Delta$ is not a perfect square), the \textbf{norm of an element} $x + y\sqrt\Delta \in K$ (where $x, y \in \Q$) is given by \cite[Ch.\,6.2]{buell}
$$N(x + y\sqrt\Delta) = (x + y\sqrt\Delta) \cdot (x - y\sqrt\Delta) = x^2 - y^2\Delta.$$


\pagebreak
\section{Ideal Correspondence}\label{s-comp-ideals}
We now seek to use the findings of \Cref{s-comp-algebraic} to develop a notion of equivalence between ideals of $\order_\Delta$ and binary quadratic forms of discriminant $\Delta$.
To do this, consider the norm of a general element of $I = l\Z + \left( m + n\frac{\Delta + \sqrt\Delta}{2} \right)\Z \trianglelefteq \order_\Delta$: 
\begin{align*}
    &\hspace{-5mm} N\left(lx + \left(m + n\frac{\Delta + \sqrt\Delta}{2}\right)y\right)
    =  N \left( \left( lx + my + n\frac{\Delta}{2}y \right) + \frac{ny}{2}\sqrt\Delta \right) \\
    &= \left( lx + my + n\frac{\Delta}{2}y \right)^2 - \left(\frac{ny}{2} \right)^2 \Delta \\
    &= l^2x^2 + m^2y^2 + n^2\frac{\Delta^2}{4}y^2 + 2lmxy + ln\Delta xy + mn\Delta y^2 - n^2\frac{\Delta}{4}y^2 \\
    &= l^2x^2 + l(2m + n\Delta)xy + \left(m^2 + mn\Delta + n^2\frac{\Delta^2 - \Delta}{4}\right)y^2. \stepcounter{equation}\tag{\theequation}\label{eq-ideal-norm-form}
\end{align*}
This is a binary quadratic form.
The discriminant of this form is given by
\begin{align*}
    &\hspace{-5mm} [l(2m + n\Delta)]^2 - 4l^2\left[ m^2 + mn\Delta + n^2\frac{\Delta^2 - \Delta}{4} \right] \\
    &= l^2(4m^2 + 4mn\Delta + n^2\Delta^2) - l^2(4m^2 + 4mn\Delta + n^2\Delta^2 - n^2\Delta) \\
    &= (ln)^2\Delta.
\end{align*}
By \Cref{lemma-discriminant-under-multiple}, if we can divide the form by $ln = N(I)$, we will arrive at a form of discriminant $\Delta$.
That is,
$$\frac{N \left( lx + \left( m + n\frac{\Delta + \sqrt\Delta}{2} \right) y \right)}{N \left( l\Z + \left( m + n\frac{\Delta + \sqrt\Delta}{2} \right) \Z \right)} = \frac{l}{n}x^2 + \left(\frac{2m}{n} + \Delta\right)xy + \frac{1}{l}\left(\frac{m^2}{n} + m\Delta + n\frac{\Delta^2 - \Delta}{4}\right)y^2.$$
The divisibility conditions of \Cref{lemma-lattice-divisibility-conditions} show each coefficient to be an integer, hence this is a binary quadratic form of discriminant $\Delta$.
We can therefore define a map $\Theta$ that takes an ideal of $\order_\Delta$ to a corresponding form of discriminant $\Delta$:
$$\Theta : l\Z + \left( m + n\frac{\Delta + \sqrt\Delta}{2} \right) \Z \longmapsto \left( \frac{l}{n}, \frac{2m}{n} + \Delta, \frac{1}{l}\left(\frac{m^2}{n} + m\Delta + n\frac{\Delta^2 - \Delta}{4}\right) \right).$$

Recall that for a given ideal, $l$ and $n$ are uniquely determined, but $m$ is only determined up to adding multiples of $l$.
We must therefore consider what effect making different choices for $m$ has on the output of this map in order to ensure it is well defined.
For our ideal $I = l\Z + \left( m + n\frac{\Delta + \sqrt\Delta}{2} \right) \Z$, we can also write 
$$I = l\Z + \left( (m + kl) + n\frac{\Delta + \sqrt\Delta}{2} \right) \Z$$
where $k \in \Z$. 
The map $\Theta$ sends this ideal to
\begin{align*}
    &\hspace{-5mm} \left( \frac{l}{n}, \frac{2(m + kl)}{n} + \Delta, \frac{1}{l}\left(\frac{(m + kl)^2}{n} + (m + kl)\Delta + n\frac{\Delta^2 - \Delta}{4}\right) \right) \\
    &= \left( \frac{l}{n}, \left[ \frac{2m}{n} + \Delta \right] + 2k\frac{l}{n}, \left[ \frac{1}{l}\left(\frac{m^2}{n} + m\Delta + n\frac{\Delta^2 - \Delta}{4}\right) \right] + k^2\frac{l}{n} + k(\frac{2m}{n} + \Delta) \right).
\end{align*}
This is the same result as a $T^k$ transformation on the previous form.
Hence we see that different choices of $m$ map to properly equivalent forms, so the map is well defined as a map into equivalence classes of forms.

To establish a bijection between ideals and forms, this map must have an inverse.
Note that an integer multiple of an ideal, $kI$, is sent to the same form as $I$, since $k$ distributes over $l, m, n$, which will then be canceled when computing the form.
So for an inverse map to be well defined, we must consider equivalence classes of ideals which associate ideals that vary be integer multiples to the same class.
In fact, we will see that the correct choice of equivalence classes is actually to associate ideals which vary by multiples of elements of $\Q(\sqrt\Delta)$ to the same class.
Now, the inverse map must take the form
\begin{align*}
    \Theta^{-1} : (a, b, c) &= \left( \frac{l}{n}, \frac{2m}{n} + \Delta, \frac{1}{l}\left(\frac{m^2}{n} + m\Delta + n\frac{\Delta^2 - \Delta}{4}\right) \right) \\
    &\mapsto l\Z + \left( m + n\frac{\Delta + \sqrt\Delta}{2} \right) \Z 
\end{align*}
where $I = l\Z + \left( m + n \frac{\Delta + \sqrt\Delta}{2} \right) \Z$ is an ideal of $\order_\Delta$ - we showed the necesarry conditions for this in \Cref{lemma-lattice-divisibility-conditions}.
Choosing $n = 1$ immediately satisfies the conditions $n \divides m$ and $n \divides l$, and  forces $l = a$ and $m = \frac{b - \Delta}{2}$ (which is always an integer, as $b \equiv \Delta \pmod{4}$).
Finally, $c = \frac{1}{l}\left(\frac{m^2}{n} + m\Delta + n\frac{\Delta^2 - \Delta}{4}\right)$ is an integer, which satisfies the final divisibility condition.
This gives the inverse map
$$(a, b, c) \mapsto a\Z + \left( \frac{b - \Delta}{2} + \frac{\Delta + \sqrt\Delta}{2} \right) \Z = a\Z + \left( \frac{b + \sqrt\Delta}{2} \right) \Z.$$

\todo{show that equivalent forms map to equivalent ideals}


\pagebreak
\section{Composition Formula}\label{s-comp-formula}
\todo{prove that equivalence classes of ideals respect multiplication}

Having seen that forms of the same discriminant correspond to ideals of the same ring, 
we now seek to find a general formula for an operation corresponding to ideal multipication.
Let us begin by considering forms $f = (a, b, c)$ and $g = (a', b', c')$, both of discriminant $\Delta$.
These correspond to ideals of $\mathcal{O}_\Delta$
$$I_f = a\Z + \frac{b + \sqrt\Delta}{2}\Z, \hspace{5mm} I_g = a'\Z + \frac{b' + \sqrt\Delta}{2}\Z.$$
The product of these ideals is given by
$$I = I_f \cdot I_g = aa'\Z + a\frac{b' + \sqrt\Delta}{2}\Z + a'\frac{b + \sqrt\Delta}{2}\Z + \frac{bb' + \Delta + (b + b')\sqrt\Delta}{4}\Z.$$
By \Cref{lemma-ideal-lattice}, we can express this product in the form $l\Z + \left(m + n\frac{\Delta + \sqrt\Delta}{2}\right)\Z$. 
We know that $n$ is the minimal positive coefficient of $m + n\frac{\Delta + \sqrt\Delta}{2}$ among elements of $I$, which can also be thought of as the coefficient of $\frac{\sqrt\Delta}{2}$.
So we must have
$$n = \gcd\left(a, a', \frac{b+b'}{2}\right).$$
Hence, we can choose $u, v, w \in \Z$ such that
$$au + a'v + \frac{b+b'}{2}w = n.$$
Which allows us to find an element of $I$ which satisfies this minimal condition:
\begin{align*}
    &\hspace{-5mm} a\frac{b' + \sqrt\Delta}{2}u + a'\frac{b + \sqrt\Delta}{2}v + \frac{bb' + \Delta + (b + b')\sqrt\Delta}{4}w \\
    &= \frac{ab'}{2}u + \frac{a'b}{2}v + \frac{bb' + \Delta}{4}w + n\frac{\sqrt\Delta}{2}
\end{align*}
\todo{find m closed form}

To find $l$, we use the fact that the norm of ideals is multiplicative \todo{cite source}.
The norm of $I_f$ is $a$, and the norm of $I_g$ is $a'$, hence, the norm of $I$ must be $aa'$.
But by \Cref{lemma-ideal-norm}, we have
\begin{gather*}
    aa' = N(I) = ln \\
    \implies l = \frac{aa'}{n}
\end{gather*}


\pagebreak
\section{Identity and Inverses}\label{s-comp-idinv}
\red{to do}


%===============================================================================================================================%
%                                                           CHAPTER 3                                                           %
%===============================================================================================================================%

\chapter{Factoring}\label{s-factor}
A key practical application of number theory is the factoring algorithms that its theory allows.
Factoring is a deceptively difficult task to perform efficiently.
As a baseline for comparison, let us consider the most primitive factoring algorithm - trial division.

\begin{algorithm}
    \caption{Factoring by trial division}
    \DontPrintSemicolon
    \KwData{$n \in \Z$}
    $x \gets 2$\;
    \While {$x \leq \sqrt{n}$}{
        \eIf {$n \equiv 0 \pmod{x}$}{
            \Return{$x$}
        }{
            $x \gets x + 1$\;
        }
    }
\end{algorithm}
\noindent
We will use big-$O$ notation to analyse the complexity of factoring algorithms.
Trial division has a complexity of $O(\sqrt{n})$.
If one considers the complexity not as a function of the size of the number being factored, but instead by the length of its binary representation,
we have that trial division scales exponentially with the length of the input that we are trying to factor.
For this reason, such an algorithm will sometimes be called \textbf{exponential} in some contexts.


\section{Ambiguous forms}\label{s-factor-amb}
The key insight which bridges the theory of binary quadratic forms into the realm of factoring comes from considering
forms of order 2 - that is, forms which are their own inverses. 
Such forms are called \textbf{ambiguous forms}, since they are properly equivalent to the same forms that they are improperly
equivalent to, so it is ``ambiguous'' which proper equivalence class they belong to.

For a reduced form to be properly equivalent to its inverse, we have
$$(a, b, c) \equiv (a, -b, c).$$
When is this possible?
The clearest case where the equility holds is when $b = 0$, but there are other possibilities.
If $b \neq 0$, the only way for the equivalence to hold is if $(a, -b, c)$ is not fully reduced.
Since $(a, b, c)$ is reduced, we know that $|b| \leq a \leq c$, 
so the only way for $(a, -b, c)$ not to be reduced is if $-b < 0$ and $a = |-b|$ or $a = c$.
Since $(a, b, c)$ is reduced, either of the latter conditions would imply $b > 0$, so we can now write the three cases
\begin{align*}
    &\textbf{(1)} \hspace{2mm} b = 0 \implies (a, 0, c) \\
    &\textbf{(2)} \hspace{2mm} a = b \implies (a, a, c) \\
    &\textbf{(3)} \hspace{2mm} a = c \implies (a, b, a).
\end{align*}
An ambiguous form must necessarily meet one of these conditions, and it is easy to show that each consition is sufficient:
\begin{align*}
    \textbf{(1)} \hspace{2mm}
    (a, 0, c) \cdot (a, 0, c) &= (a, 0, c) \cdot (a, -0, c) \\
    &= id_\Delta
\end{align*}
\begin{align*}
    \textbf{(2)} \hspace{2mm}
    (a, a, c) \cdot (a, a, c) &= (a, a, c) \cdot (a, -a, a - a + c) \rnote{\text{(using $T^{-1}$)}} \\
    &= (a, a, c) \cdot (a, -a, c) \\
    &= id_\Delta
\end{align*}
\begin{align*}
    \textbf{(3)} \hspace{2mm}
    (a, b, a) \cdot (a, b, a) &= (a, b, a) \cdot (a, -b, a) \rnote{\text{(using $S$)}} \\
    &= id_\Delta
\end{align*}
so these cases are necessary and sufficient to classify a form as ambiguous.
Notice that each case reduces the three variables $a, b, c$ to only two - 
this is of particular note when one considers how each case affects the formula $\Delta = b^2 - 4ac$
\begin{gather*}
    \textbf{(1)} \hspace{2mm} \Delta = 0 - 4ac = -4ac \\
    \textbf{(2)} \hspace{2mm} \Delta = a^2 - 4ac = a(a-4c) \\
    \textbf{(3)} \hspace{2mm} \Delta = b^2 - 4a^2 = (b + 2a)(b - 2a).
\end{gather*}
Notice that every case gives a factorisation of the discriminant.
Therefore, if we want to find a factorisiation of some number $D$, we can do this by finding amabiguous forms of discriminant $\Delta = D$
(or $\Delta = 4D$ if $D$ is congruent to 2 or 3 modulo 4).


%===============================================================================================================================%
%                                                           REFERENCES                                                          %
%===============================================================================================================================%

\pagebreak
\bibliographystyle{amsalpha}
\bibliography{references}
\addcontentsline{toc}{chapter}{References}
\end{document}